\documentclass[11pt]{article}
\usepackage[margin=1in]{geometry}
\usepackage{amsmath,amssymb}
\usepackage{booktabs}
\usepackage{enumitem}
\usepackage{xcolor}
\usepackage{graphicx}
\usepackage[hidelinks]{hyperref}
\setlist{itemsep=2pt,topsep=4pt}

\definecolor{accent}{HTML}{B4540A}
\newcommand{\key}[1]{\par\smallskip\noindent\fbox{\parbox{\dimexpr\linewidth-2\fboxsep-2\fboxrule}{\textbf{Key point.} #1}}\par\smallskip}

\IfFileExists{baselines/results/tables.tex}{% generated by baselines/make_tables.py from raw.json (3 runs, 3 splits)
\newcommand{\nruns}{3}
\newcommand{\nsplits}{3}
\newcommand{\nodetable}{\begin{tabular}{llcccccc}
\toprule
Model & Attack & Acc.\,$\uparrow$ & NLL\,$\downarrow$ & ECE\,$\downarrow$ & Brier\,$\downarrow$ & Acc.\ touched & Acc.\ untouched \\
\midrule
MLP (no edges) & -- & 0.738\,$\pm$\,0.010 & 0.781\,$\pm$\,0.018 & 0.021\,$\pm$\,0.011 & 0.371\,$\pm$\,0.008 & -- & -- \\
\midrule
GCN & none & 0.794\,$\pm$\,0.005 & 0.724\,$\pm$\,0.015 & 0.043\,$\pm$\,0.006 & 0.310\,$\pm$\,0.009 & -- & 0.794\,$\pm$\,0.005 \\
\bottomrule
\end{tabular}}
\newcommand{\aurocTable}{\begin{tabular}{lcccccc}
\toprule
& \multicolumn{2}{c}{Text only} & \multicolumn{2}{c}{Structure only} & \multicolumn{2}{c}{Naive mix} \\
\cmidrule(lr){2-3}\cmidrule(lr){4-5}\cmidrule(lr){6-7}
Attack & SBERT & GloVe & Jaccard & GCN JS & SBERT+GCN & SBERT+Jacc. \\
\midrule
\bottomrule
\end{tabular}}
\newcommand{\auprcTable}{\begin{tabular}{lcccccc}
\toprule
& \multicolumn{2}{c}{Text only} & \multicolumn{2}{c}{Structure only} & \multicolumn{2}{c}{Naive mix} \\
\cmidrule(lr){2-3}\cmidrule(lr){4-5}\cmidrule(lr){6-7}
Attack & SBERT & GloVe & Jaccard & GCN JS & SBERT+GCN & SBERT+Jacc. \\
\midrule
\bottomrule
\end{tabular}}
}{%
  \newcommand{\nodetable}{\emph{(run \texttt{python baselines/make\_tables.py})}}%
  \newcommand{\aurocTable}{\nodetable}\newcommand{\auprcTable}{\nodetable}%
  \newcommand{\nruns}{0}\newcommand{\nsplits}{0}}

\title{Which Edges Can a GNN Trust?\\[4pt]
\large Baseline results on WikiCS}
\author{Luis Mendez \quad CSE5825 Bayesian Machine Learning}
\date{October 2, 2026}

\begin{document}
\maketitle

%================================================================
\section{Summary}
%SUMMARY%

%================================================================
\section{Setup}

\paragraph{Data.} WikiCS \cite{wikics}: 11{,}701 Wikipedia articles on computer science,
215{,}603 undirected hyperlink edges, and 10 topic labels. Each node has its raw article text
(title plus tokens) and a 300-dimensional averaged GloVe feature vector. Clean edge homophily
(the share of edges joining two nodes with the same label) is 0.654. I use the first
\nsplits{} of the 20 standard train/validation splits (580 training nodes each) and the
standard test set of 5{,}847 nodes. The data was loaded from the original JSON release.

\paragraph{Attack.} Fake edges are injected between nodes whose texts look related but whose
labels differ. For every node I take its 60 nearest neighbours by SBERT cosine similarity
(\texttt{all-MiniLM-L6-v2} on title plus the first 256 tokens) and keep the non-adjacent pairs
with different labels. This gives 201{,}981 candidate pairs. An attack at budget $b$ samples
$b \cdot 215{,}603$ of them uniformly, for $b \in \{5, 10, 20, 40\}\%$. As a contrast, a
\emph{random} attack at $b = 20\%$ samples non-edges uniformly, ignoring text and labels.
Because I inject the fake edges myself, every edge has a known ground-truth label (real or
fake).

\paragraph{Models.} A 2-layer GCN \cite{gcn} and a 2-layer MLP that ignores edges, both with
64 hidden units, dropout 0.5, Adam (learning rate 0.01, weight decay $5\times10^{-4}$), up to
300 epochs, and early stopping on validation loss (patience 50). Node features are the GloVe
vectors. The GCN at budget 0\% is also the \emph{oracle defense}: it is what any method that
removed exactly the fake edges would achieve.

\paragraph{Node metrics.} Test accuracy, negative log-likelihood (NLL), expected calibration
error (ECE, 15 bins) and Brier score. I also split test accuracy by whether a node is an
endpoint of at least one fake edge (\emph{touched}) or not.

\paragraph{Edge-detection baselines.} Each method gives every edge of the attacked graph a
suspicion score, higher meaning more likely fake. Scores are compared with the ground truth by
AUROC and AUPRC, with fake edges as the positive class.
\begin{itemize}
  \item \textbf{Text only:} $1 - \cos$ of the two endpoints' SBERT embeddings, and the same for
        their GloVe features. These stand in for the text prior until the LLM prior is added.
  \item \textbf{Structure only:} $1 -$ Jaccard overlap of the endpoints' neighbourhoods, and
        the Jensen--Shannon divergence between the GCN's predicted class distributions at the
        two endpoints (the GCN trained on the attacked graph).
  \item \textbf{Naive mix:} the unweighted average of the rank-normalised text and structure
        scores. This is the ``hand-made rule'' that the Bayesian model should beat.
\end{itemize}

%================================================================
\section{Results}

\subsection{Node classification}

\begin{table}[h]
\centering\small
\resizebox{\linewidth}{!}{\nodetable}
\caption{Test-set node classification on WikiCS, mean\,$\pm$\,std over \nsplits{} splits.
``Touched'' nodes are endpoints of at least one fake edge.}
\label{tab:node}
\end{table}

%NODE-DISCUSSION%

\subsection{Fake-edge detection}

\begin{table}[h]
\centering\small
\aurocTable
\caption{Fake-edge detection AUROC (fake = positive; 0.5 = chance, below 0.5 = ranks fakes
as \emph{more} trustworthy). Mean over \nsplits{} splits; best per row in bold.}
\label{tab:auroc}
\end{table}

\begin{table}[h]
\centering\small
\auprcTable
\caption{Fake-edge detection AUPRC. The chance level equals the share of fake edges:
0.048, 0.091, 0.167 and 0.286 at budgets 5, 10, 20 and 40\%.}
\label{tab:auprc}
\end{table}

%DETECT-DISCUSSION%

%================================================================
\section{What this means for the hypotheses}
%HYPOTHESES%

%================================================================
\section{Limitations and next steps}
%NEXT%

\begin{thebibliography}{9}
\bibitem{wikics} P.~Mernyei and C.~Cangea. Wiki-CS: A Wikipedia-based benchmark for graph
  neural networks. \emph{ICML Workshop on Graph Representation Learning}, 2020.
\bibitem{gcn} T.~N.~Kipf and M.~Welling. Semi-supervised classification with graph
  convolutional networks. \emph{ICLR}, 2017.
\end{thebibliography}

\end{document}
